\chapter*{Abstract}

This thesis presents a formal, declarative model for questionnaire specification that fundamentally differs from traditional flow-based approaches. Rather than viewing questionnaires as sequential execution paths, we formalize them as collections of items with preconditions (Boolean predicates determining when questions appear) and postconditions (constraints on acceptable responses). This declarative paradigm, inspired by Hoare Logic, treats the questionnaire as a set of logical relationships where the actual flow emerges dynamically based on respondent answers. Our static analysis investigates all possible paths simultaneously through constraint satisfaction, reasoning about relationships between immutable states rather than simulating execution through mutable values.

Our implementation leverages the Z3 SMT solver to provide mathematical guarantees about questionnaire correctness. Despite questionnaire consistency checking being NP-complete (Willenborg, 1995), modern SMT solvers handle real-world questionnaires effectively. This enables validation of questionnaires containing hundreds of interconnected questions with complex conditional dependencies---a level of complexity that exceeds human comprehension and often makes traditional validation approaches computationally intractable.
